\section{Evaluasi}
\label{sec:evaluasi-bab7}

\subsection{Soal Pilihan Ganda}

\begin{enumerate}
    \item Berkas keluaran yang dihasilkan langsung oleh TASM adalah
          \dots
          \begin{enumerate}
              \item \file{.ASM} \item \file{.OBJ}
              \item \file{.EXE} \item \file{.COM}
          \end{enumerate}

    \item Perangkat yang menggabungkan berkas objek menjadi berkas
          yang dapat dijalankan adalah \dots
          \begin{enumerate}
              \item editor \item TASM
              \item TLINK \item DEBUG
          \end{enumerate}

    \item Direktif yang wajib dipakai pada program \file{.COM}
          adalah \dots
          \begin{enumerate}
              \item \dir{.MODEL SMALL} \item \dir{ORG 100h}
              \item \dir{.STACK 100h} \item \dir{ASSUME CS:CODE}
          \end{enumerate}

    \item Direktif yang menyediakan ruang dua byte untuk satu nilai
          adalah \dots
          \begin{enumerate}
              \item \dir{DB} \item \dir{DW}
              \item \dir{DD} \item \dir{EQU}
          \end{enumerate}

    \item Inisialisasi \reg{DS} diperlukan pada program \file{.EXE}
          karena \dots
          \begin{enumerate}
              \item DOS tidak menetapkan \reg{DS} ke segmen data
                    program
              \item \reg{DS} hanya boleh diisi oleh TLINK
              \item program \file{.EXE} tidak memiliki segmen data
              \item \reg{DS} harus sama dengan \reg{CS}
          \end{enumerate}
\end{enumerate}

\subsection{Soal Uraian}

\begin{enumerate}
    \item Jelaskan rantai proses dari berkas \file{.ASM} sampai
          program berjalan. Sebutkan peran editor, TASM, TLINK, dan
          \eng{loader}.

    \item Sebutkan empat perbedaan program \file{.COM} dan
          \file{.EXE}. Jelaskan mengapa program \file{.COM} tidak
          memerlukan inisialisasi \reg{DS}, sedangkan program
          \file{.EXE} memerlukannya.

    \item Jelaskan fungsi \dir{DB}, \dir{DW}, \dir{DD}, \dir{EQU},
          dan \dir{DUP}. Berikan satu contoh pemakaian untuk
          masing-masing direktif.

    \item Perhatikan program berikut.
          \begin{lstlisting}[language=MASM, numbers=none]
.MODEL SMALL
.STACK 100h
.DATA
    nilai DW 1234h
.CODE
START:
    MOV AX, @DATA
    MOV DS, AX
    MOV AX, nilai
    MOV AH, 4Ch
    INT 21h
END START
          \end{lstlisting}
          Jelaskan fungsi setiap direktif dan label pada program
          tersebut. Jelaskan pula apa yang terjadi bila
          \lstasm|MOV DS, AX| dihapus.
\end{enumerate}

\subsection{Kisi-kisi Penilaian}

\begin{table}[H]
    \centering
    \caption{Kisi-kisi evaluasi Bab~\ref{chap:memulai-assembly}.}
    \label{tab:kisi-bab7}
    \small
    \begin{tabular}{@{}p{1.5cm}p{5.0cm}p{3.0cm}c@{}}
        \toprule
        \textbf{Sub-CPMK} & \textbf{Indikator yang diukur} &
        \textbf{Bentuk soal} & \textbf{Bobot} \\
        \midrule
        2.1 & Menjelaskan rantai proses dan peran TASM/TLINK &
        Pilihan ganda 1, 2; Uraian 1 & 30\% \\
        \addlinespace
        2.1 & Membedakan program COM dan EXE &
        Pilihan ganda 3, 5; Uraian 2 & 30\% \\
        \addlinespace
        2.1 & Menjelaskan struktur program dan direktif &
        Pilihan ganda 4; Uraian 3, 4 & 40\% \\
        \bottomrule
    \end{tabular}
\end{table}

\begin{checklistbox}
\vspace{1.1em}
\textbf{Checklist Pencapaian Kompetensi
Bab~\ref{chap:memulai-assembly}}

Seorang mahasiswa dinyatakan kompeten pada Sub-CPMK 2.1 apabila:
\begin{itemize}
    \item[\checked] dapat menjelaskan alur
          \file{.ASM} $\rightarrow$ \file{.OBJ} $\rightarrow$
          \file{.EXE}/\file{.COM};
    \item[\checked] dapat menjelaskan peran TASM dan TLINK;
    \item[\checked] dapat membedakan program \file{.COM} dan
          \file{.EXE};
    \item[\checked] dapat menjelaskan fungsi \dir{.MODEL},
          \dir{.STACK}, \dir{.DATA}, \dir{.CODE}, \dir{ORG 100h},
          dan \dir{END};
    \item[\checked] dapat menulis label dan komentar dengan benar;
    \item[\checked] dapat menggunakan direktif data \dir{DB},
          \dir{DW}, \dir{DD}, \dir{EQU}, dan \dir{DUP};
    \item[\checked] dapat meng-assemble, men-link, dan menjalankan
          program sederhana.
\end{itemize}
\end{checklistbox}
